\section*{Limitations}

\textbf{Process-metric scope.}
C1--C3 are reproducible proxies grounded in verifier scores and execution signals rather than exhaustive definitions of the underlying research capabilities.
They do not capture aspects that are not reliably visible in the trajectory, such as the semantic quality of an unrealized idea or the agent's latent reasoning.
Their evidential strength also depends on the events observed during a run: in particular, C3 cannot meaningfully measure recovery when a trajectory contains few or no regressions.
A short or nearly monotone trajectory may therefore receive a high Feedback Control score without demonstrating recovery from repeated setbacks.
The three metrics should be interpreted together with the behavioral diagnostics and trajectory evidence rather than as standalone measures of general research ability.

\textbf{Controlled-experiment dependence.}
Our self-improvement estimates depend on the interventions used to isolate experience, including the intra-task erasure point, the selected source--target pairs, and the representation of transferred experience.
These controls support causal comparisons within the evaluated settings, but they do not cover every way an agent might accumulate, retrieve, revise, or forget experience over a longer deployment.
Different memory systems or task sequences may therefore produce different estimates of experience-driven improvement.

\textbf{Benchmark and harness dependence.}
Our conclusions are based on the task distribution, resource budgets, verifiers, and execution environment of AutoLab, with a common harness used for the main cross-model comparison.
Although the harness experiments show which findings are stable under several alternative scaffolds, they do not exhaust the space of prompts, tools, context-management policies, or model--harness combinations.
Absolute scores and some relative rankings may change under other research domains or system configurations.

\textbf{Cost comparability.}
Estimated inference cost depends on provider pricing, token accounting, and serving configurations, all of which can change over time or differ across deployments.
Wall-clock use is also affected by task-specific budgets and infrastructure conditions.
Cost results are therefore most reliable for comparisons under our controlled setup and should not be interpreted as universal deployment prices.
